\documentclass[10pt]{article}

% ---------- Document Identity (update these only) ----------
\newcommand{\DocStage}{}
\newcommand{\DocTitle}{Your Road to Tier One SOF}
\newcommand{\ProgramName}{182-Week Civilian-to-Selection Readiness Pipeline}
\newcommand{\ProgramSubtitle}{}

% ---------- Page ----------
\usepackage[legalpaper,left=0.5in,right=0.5in,top=0.6in,bottom=0.45in]{geometry}

% ---------- Typography ----------
\usepackage[T1]{fontenc}
\usepackage[utf8]{inputenc}
\usepackage{libertinus}
\usepackage{microtype}
\usepackage{parskip}
\usepackage{ragged2e}
\usepackage{textcomp}
\AtBeginDocument{\color{black}}
\AtBeginDocument{\pagecolor{white}}
\AtBeginDocument{\justifying}

% ---------- Landscape pages ----------
\usepackage{pdflscape}

% ---------- Color ----------
\usepackage[table]{xcolor}
\definecolor{StageBlue}{HTML}{1F4E79}
\definecolor{StageBlueDark}{HTML}{163A5A}
\definecolor{LightGray}{HTML}{F2F4F7}
\definecolor{MidGray}{HTML}{D9DEE5}
\definecolor{RuleGray}{HTML}{C7CED8}
\definecolor{HeaderInk}{HTML}{000000}
\definecolor{Ink}{HTML}{000000}

% ---------- Utility ----------
\usepackage{enumitem}
\sloppy
\setlength{\emergencystretch}{2em}

% ---------- Boxes / UI ----------
\usepackage[most]{tcolorbox}

\tcbset{
  charterTitle/.style={
    enhanced,
    colback=StageBlue!4,
    colframe=StageBlue,
    boxrule=1.25pt,
    arc=3pt,
    left=16pt,right=16pt,top=14pt,bottom=14pt,
    borderline north={3.2pt}{0pt}{StageBlue},
    borderline west={4pt}{0pt}{StageBlue},
    borderline south={1.25pt}{0pt}{StageBlue},
    drop shadow={black!25!white},
  },
  charterRule/.style={
    enhanced,
    colback=LightGray,
    colframe=RuleGray,
    boxrule=0.85pt,
    arc=3pt,
    left=12pt,right=12pt,top=10pt,bottom=10pt,
    borderline west={3pt}{0pt}{StageBlue},
  },
  charterBadge/.style={
    enhanced,
    colback=StageBlue,
    coltext=white,
    colframe=StageBlueDark,
    boxrule=0.6pt,
    arc=2pt,
    left=6pt,right=6pt,top=3pt,bottom=3pt,
  }
}
\newtcbox{\Badge}{nobeforeafter,charterBadge}

% ---------- Lists ----------
\setlist[itemize]{leftmargin=1.5em, itemsep=0.25em, topsep=0.25em}
\setlist[enumerate]{leftmargin=1.8em, itemsep=0.25em, topsep=0.25em}

% ---------- TOC ----------
\usepackage{tocloft}
\setcounter{tocdepth}{3}
\setlength{\cftbeforesecskip}{3pt}
\setlength{\cftbeforesubsecskip}{2pt}
\setlength{\cftbeforesubsubsecskip}{0pt}
\renewcommand{\cftsecfont}{\bfseries}
\renewcommand{\cftsecpagefont}{\bfseries}
\renewcommand{\cftsubsecfont}{\normalfont}
\renewcommand{\cftsubsecpagefont}{\normalfont}
\renewcommand{\cftsubsubsecfont}{\normalfont}
\renewcommand{\cftsubsubsecpagefont}{\normalfont}
\setlength{\cftsecindent}{0pt}
\setlength{\cftsubsecindent}{1.25em}
\setlength{\cftsubsubsecindent}{2.8em}
\setlength{\cftsecnumwidth}{2.1em}
\setlength{\cftsubsecnumwidth}{2.8em}
\setlength{\cftsubsubsecnumwidth}{3.6em}

% Remove the built-in TOC heading so only the custom heading is shown
\makeatletter
\renewcommand{\tableofcontents}{\@starttoc{toc}}
\makeatother

% ---------- Header/Footer: page number only ----------
\usepackage{fancyhdr}
\pagestyle{fancy}
\fancyhf{}
\renewcommand{\headrulewidth}{0pt}
\renewcommand{\footrulewidth}{0pt}
\fancyfoot[C]{\thepage}

% ---------- Section formatting ----------
\usepackage{titlesec}

\titleformat{\section}
  {\Large\bfseries\color{Ink}}
  {\thesection}{0.6em}{}
  [\vspace{0.25em}\textcolor{StageBlue}{\rule{\linewidth}{0.8pt}}]

\titleformat{\subsection}
  {\large\bfseries\color{Ink}}
  {\thesubsection}{0.6em}{}

\titleformat{\subsubsection}
  {\normalsize\bfseries\color{Ink}}
  {\thesubsubsection}{0.6em}{}

\titlespacing*{\section}{0pt}{0.95em}{0.35em}
\titlespacing*{\subsection}{0pt}{0.65em}{0.25em}
\titlespacing*{\subsubsection}{0pt}{0.55em}{0.2em}

% ---------- Table engine ----------
\usepackage{tabularray}
\UseTblrLibrary{booktabs}

% ---------- tabularray: GLOBAL table treatment ----------
\SetTblrInner{
  cells = {fg=black, bg=white, valign=t},
}

% ---------- Header helpers ----------
\newcommand{\Hdr}[1]{\shortstack[c]{\strut #1\strut}}

% ---------- Lines ----------
\newcommand{\TitleLine}{\textcolor{StageBlue}{\rule{0.98\linewidth}{0.95pt}}}
\newcommand{\SubTitleLine}{\textcolor{RuleGray}{\rule{0.98\linewidth}{0.65pt}}}

% ---------- Continued on next page ----------
\newcommand{\ContinuedOnNextPageInline}{%
  \par
  \vfill
  \noindent\textcolor{RuleGray}{\rule{\linewidth}{1pt}}\vspace{-2pt}
  \noindent\hfill\textit{Continued on next page}\par\vspace{-10pt}
  \noindent\textcolor{RuleGray}{\rule{\linewidth}{1pt}}
  \clearpage
}

% ---------- PDF hyperlinks ----------
\usepackage[hidelinks]{hyperref}
\hypersetup{
  colorlinks=false,
  pdfborder={0 0 0},
  linktoc=all,
  pdftitle={\DocTitle},
  pdfauthor={\ProgramName}
}

% =========================================================
\begin{document}

\begin{tcolorbox}
\begin{center}
Stage 1 — Charter, Governance, and Templates\\[0.35em]
SOF Physical Development Transformation Program\\[0.25em]
182-Week Civilian-to-Selection Readiness Pipeline\\[0.65em]
\TitleLine
\end{center}
\end{tcolorbox}

\vspace{0.35em}

\section*{Stage 1 — Charter, Governance, and Templates}
\phantomsection
\addcontentsline{toc}{section}{Stage 1 — Charter, Governance, and Templates}

\subsection*{Introduction}
\phantomsection
\addcontentsline{toc}{subsection}{Introduction}

\subsubsection*{1) Purpose}

Stage 1 exists to function as the governance authority layer for the full 182-week civilian-to-selection readiness pipeline, Phase 00 through Phase 13. Stage 1 locks the operational definitions, compliance doctrine, evidence posture, and change-control mechanics that prevent execution drift, prevent unsafe escalation, and preserve audit-ready comparability over a multi-year horizon.

Stage 1 \textbf{MUST} achieve the following:

\begin{itemize}
\item Stabilize decision language. Stage 1 locks canonical definitions and stable terminology so later artifacts cannot silently change meaning to justify convenience, speed, or preference. If definitions drift, evidence becomes non-comparable, and gate decisions become subjective.

\item Convert Stage 0 constants into enforceable governance controls. Stage 0 defines what is fixed; Stage 1 defines how those fixed rules are applied operationally in day-to-day decision-making, including what counts, what does not count, and what happens when compliance fails.

\item Make execution auditable. Stage 1 defines the minimum structures and required schemas so a third-party coach can audit execution and reach the same decisions using the same admissible evidence without needing oral context, personal familiarity, or informal “coach intuition.”

\item Make risk an owned control problem. Stage 1 defines how risks are tracked, detected early via triggers or signals, mitigated through explicit controls, and owned by a single accountable Owner. Risk posture is operational and enforceable, not rhetorical.

\item Make dependencies explicit and validated. Stage 1 ensures prerequisites are not “assumed until failure.” Dependencies have validation methods, due-by timing, Owners, and explicit fail actions that enforce substitution, delay, Hold, Regress, or Medical Hold as required.

\item Prevent mid-pipeline moving targets. Stage 1 establishes freeze-window and patch discipline so evidence remains comparable and so changes do not corrupt gates, invalidate tests, or create hidden contradictions between stages.
\end{itemize}

Stage 1 does not prescribe day-to-day training. Any downstream artifact is non-deployable if it violates Stage 0 or Stage 1.

\subsubsection*{2) Scope}

\paragraph*{2.1 Inclusions}

Stage 1 controls compliance and governance across Phase 00 through Phase 13 by deliverable.

\begin{itemize}
\item 1.A Project Charter

Locks mission framing, scope boundaries, non-negotiables, definition-of-success posture aligned to Stage 0 targets, decision hierarchy, roles and decision rights, enforcement posture, and charter-level change-control posture. 1.A \textbf{MUST} define what is fixed, what is allowed to vary, and who has stop authority, gate authority, and patch authority.

\item 1.B Key Risks and Mitigations

Defines the owned risk register with triggers or signals, impacts, likelihood, mitigations, and Owners. 1.B \textbf{MUST} convert predictable failure modes into a detection-and-response system that constrains gates and prevents unsafe escalation.

\item 1.C Assumptions and Dependencies

Defines prerequisites as validated dependencies with rationale, validation method, due-by timing, Owner assignment, and explicit fail actions. 1.C \textbf{MUST} enforce substitute, delay, Hold, Regress, or Medical Hold posture when dependencies are unmet, and \textbf{MUST} prevent “pretend ready” assumptions.

\item 1.D Milestone Plan and Output Order

Locks stage build sequence, prerequisite posture, and internal order controls that prevent out-of-order artifact generation. 1.D \textbf{MUST} implement no-skip rules and define how out-of-order work is handled: non-deployable until reconciled and re-versioned.

\item 1.E Standardized Templates

Locks the standardized templates and table schemas used downstream so artifacts remain complete, comparable, and crosswalk-ready. 1.E \textbf{MUST} provide mandatory schemas with stable headers and required fields.

Includes:

1.E.1 Overview Template

1.E.2 Phase Row Template

\item 1.F Governance Rules and Validity Doctrine

Locks validity tagging, gate outcomes, stop-rule integration, freeze windows, patch protocol, change-log requirements, and admissibility rules for gate decisions. 1.F \textbf{MUST} define the rules-of-evidence layer that prevents invalid sessions/tests from being used to justify advancement, and \textbf{MUST} formalize stop posture as overriding schedule and performance.
\end{itemize}

Stage 1 is the compliance authority across Phase 00–Phase 13. Stage 1 operationalizes Stage 0 constants into enforceable governance controls and \textbf{MUST} NOT redefine Stage 0.

\paragraph*{2.2 Exclusions}

Stage 1 explicitly excludes:

\begin{itemize}
\item Training prescriptions, weekly microcycles, exercise programming, intervals, sets, reps, progressions, and day-to-day schedules.
\item Test protocol details and scoring standards, which belong to Stage 2.
\item Test calendars and deload maps, which belong to Stage 3.
\item Phase playbooks and phase prescriptions, which belong to Stage 4 and Stage 5.
\item Equipment build-out and overviews, which belong to Stage 6 and Stage 7.
\item Crosswalk and QA packet, which belong to Stage 8.
\item Packaging and implementation guide, which belong to Stage 9.
\end{itemize}

\subsubsection*{3) How to Use}

Stage 1 is used as the continuous governance reference layer for build and execution. Coach Lead \textbf{MUST} treat Stage 1 as the compliance checklist applied to every downstream artifact and every gate decision.

When Stage 1 \textbf{MUST} be consulted:

\begin{itemize}
\item Program start and Phase 00 start: confirm governance readiness, including role mechanics, decision rights, stop posture, validity discipline, dependency validation posture, and minimum dataset expectations.

\item Weekly execution oversight: review compliance, evidence completeness, validity tagging distribution, risk triggers or signals, pain trends, environment constraints, and dependency status changes that could alter safe execution.

\item Mid-phase and end-of-phase gates: apply gate doctrine, confirm admissible evidence posture, refuse advancement without \textbf{VALID} evidence, and document gate outcomes with traceability to dated records.

\item Test weeks and evidence verification windows: confirm that evidence is produced under recorded conditions and that no novelty or uncontrolled changes contaminate comparability.

\item Injury events, mechanics-altering pain, red-flag symptoms, and mental red flags: execute stop posture, document the event, escalate to Medical when indicated, and constrain gates until stability and admissibility are restored.

\item Environment and logistics changes: route changes, weather shifts, facility access changes, equipment failures, or schedule disruptions that affect validity or safety. Apply substitution doctrine and documentation requirements; do not permit drift or undocumented workarounds.

\item Any proposed execution change: determine whether the change is a runtime coaching adjustment within authorized levers or a doctrine-level change requiring a patch. If doctrine changes are required, initiate patch protocol; no silent edits.
\end{itemize}

Conflict resolution posture:

\begin{itemize}
\item Stage 0 governs constants: identity, scope boundaries, operating constraints, safety boundaries, end-state envelope, and package usage doctrine.
\item Stage 1 governs compliance: validity, gate mechanics, stop-rule integration, risk and dependency posture, standardized templates, and change control.
\item Downstream artifacts \textbf{MUST} comply. If they do not comply, they are non-deployable until patched, corrected, and re-versioned. No silent edits.
\end{itemize}

Decision checklist use:

\begin{itemize}
\item Stage 1 artifacts are used as decision checklists for gate outcomes: \textbf{ADVANCE}, \textbf{HOLD}, \textbf{REGRESS}, and \textbf{MEDICAL} \textbf{HOLD}. Decisions \textbf{MUST} be documented, evidence-based, and restricted to admissible evidence.
\end{itemize}

Audit posture:

\begin{itemize}
\item Stage 1 is written so a third-party coach can interpret the rules without outside context, audit compliance, and reach the same decision given the same admissible evidence. Stage 1 therefore \textbf{MUST} remain explicit, stable in terminology, and resistant to interpretation drift.
\end{itemize}

\subsubsection*{4) Inputs and Assumptions}

\paragraph*{4.1 Inputs Binding}

Stage 1 is constrained by Stage 0 as governing inputs. Stage 1 references Stage 0 without redefining it.

Inputs Stage 1 \textbf{MUST} treat as binding include, at minimum:

\begin{itemize}
\item Program structure: 182-week pipeline, Phase 00 through Phase 13 executed in order.
\item Start-state and start-from-zero posture as the governing feasibility constraint.
\item End-state target envelope and definition-of-done posture, including repeatability and audit-grade evidence requirements.
\item Operating constraints: 6 sessions per week, fixed session structure, Cardio minimum, steps posture as a separate governance anchor, and the linkage between constraints and session validity.
\item Safety, legal, and medical boundaries: OTC-only posture for guidance, no prescription guidance, no hormone dosing/sourcing/cycling/procurement, and strict water safety governance including prohibition on unsupervised underwater or breath-hold exposure.
\item Execution and update discipline: validity-only evidence for advancement decisions, freeze windows, patch protocol, and conflict hierarchy.
\end{itemize}

\paragraph*{4.2 Assumption Ledger: Only If Needed}

\begin{itemize}
\item A low-friction daily logging method exists and is usable enough to support validity tagging, minimum dataset capture, and trend review.
\item Owners can be assigned and held accountable using the approved Owner taxonomy: Coach Lead, Trainee, Medical, Equipment Lead.
\item Stop rules can be executed immediately when triggers or signals appear, and are treated as mandatory actions rather than discretionary preferences.
\item Dependencies can be validated by explicit methods and due-by timing, and are not treated as “assumed true” until failure.
\item Change control can be applied operationally using freeze windows and patch discipline, including safety-critical patch posture when required.
\end{itemize}

\subsubsection*{5) Interfaces}

\paragraph*{5.1 Upstream Dependencies}

\begin{itemize}
\item Stage 0, 0.1 through 0.6, is authoritative and constrains Stage 1.
\item Stage 1 \textbf{MUST} NOT redefine Stage 0 constants; it operationalizes them into enforceable governance controls.
\end{itemize}

\paragraph*{5.2 Downstream Consumers}

\begin{itemize}
\item Stages 2 through Stage 12 depend on Stage 1 for stable terminology, standardized templates, validity posture, gate mechanics, stop-rule integration, and change-control discipline.
\item During execution, Stage 1 remains a continuous reference for compliance evaluation, risk response, dependency validation, gate decisions, and controlled updates.
\end{itemize}

\paragraph*{5.3 Crosswalk Hooks}

Stage 1 defines crosswalk-critical governance elements used everywhere downstream and treated as stable keys:

\begin{itemize}
\item Governance tags and outcomes:

\textbf{VALID}, \textbf{MODIFIED}, \textbf{INVALID}, ,\textbf{ADVANCE}, \textbf{HOLD}, \textbf{REGRESS}, \textbf{MEDICAL} \textbf{HOLD}

Risk flag

Dependency

Owner

Trigger or signal

\item Template structures that downstream artifacts \textbf{MUST} use:

Overview Template

Phase Row Template

\item Change-control posture:

versioning, freeze windows, patch protocol, and change log entries that preserve evidence comparability and require downstream revalidation after governance-layer changes.
\end{itemize}

\subsubsection*{6) Gate / Acceptance Criteria}

Stage 1 is complete and deployable only when all criteria below are satisfied:

\begin{itemize}
\item Deliverables 1.A through 1.F are produced.
\item No contradictions exist with Stage 0 constants. Any contradiction renders Stage 1 non-deployable until corrected and re-versioned.
\item Required tables use required columns with no renamed headers and no missing decision-critical fields; standardized templates remain stable and enforceable.
\item Validity doctrine is explicit and enforceable such that sessions and tests can be classified consistently and gate decisions are restricted to admissible evidence.
\item Change-control posture is explicit and enforceable, including freeze windows, patch protocol, and approval triggers for high-risk changes.
\item Stop rules and escalation pathways are unambiguous, apply to both training sessions and test events, and override schedule and performance.
\item A third-party coach can audit compliance, evidence quality, and artifact deployability without outside context and reach the same decision given the same admissible evidence.
\end{itemize}

\subsubsection*{7) Common Failure Modes and Controls}

\begin{itemize}
\item Governance drift, terms redefined midstream → locked terminology and standardized templates; Stage 0 and Stage 1 override hierarchy enforced; non-deployable rule applied.
\item Vague risk handling without triggers or signals → owned risk register requires observable triggers or signals, impacts, mitigations, and a single Owner.
\item Dependencies assumed until failure → dependency table requires validation method, due-by timing, Owner, and fail action; validation is enforced before escalation.
\item Out-of-order build sequence → milestone plan gates prevent downstream artifacts from being produced without prerequisites; out-of-order artifacts are non-deployable until reconciled and re-versioned.
\item Ad hoc mid-phase changes → freeze windows and patch protocol; explicit entries and validation checks required; no silent edits.
\item Unsafe escalation under motivation or time pressure → stop rules and regression triggers; conservative substitution posture; no make-up volume; safety overrides schedule and performance.
\item Invalid sessions or tests counted as compliance → validity-only evidence policy; no logs equals no evidence; no evidence equals no gate credit; gate outcome defaults to \textbf{HOLD} when evidence is insufficient.
\item Documentation collapse → minimum dataset requirements; gate decisions prohibited without auditable evidence; repeated missing data forces conservative gate posture.
\item Owner ambiguity and diffusion of responsibility → single accountable Owner required for each risk, dependency, gate action, and patch; unowned items are treated as governance defects.
\item Template drift and renamed columns → standardized schemas are mandatory; renamed headers or missing decision-critical fields render artifacts non-deployable.
\item Cherry-picked evidence and peak-day claims → repeatability posture enforced; anomalous outputs do not authorize advancement without admissible replication.
\item Quiet weakening of safety boundaries → \textbf{MUST} NOT rules enforced; boundary violations trigger documentation, conservative gate posture, and revalidation of affected artifacts.
\end{itemize}

\subsubsection*{8) Versioning Notes}

Frozen in Stage 1:

\begin{itemize}
\item governance terms and definitions used across the system
\item templates and locked table structures, including required headers and mandatory fields
\item validity discipline and evidence rules that constrain gate credit and admissibility
\item stop-rule posture and escalation doctrine aligned to Stage 0 safety boundaries
\item change-control posture, including freeze windows, patch discipline, and change log requirements
\end{itemize}

Permitted evolution is limited to clarifying language that does not change meaning, requirements, or enforcement posture. Any substantive change \textbf{MUST} be introduced through explicit change log entries under freeze-window discipline and \textbf{MUST} preserve alignment with Stage 0 constants.

Downstream artifacts \textbf{MUST} be revalidated after any governance-layer change. If revalidation is not performed, affected downstream artifacts are non-deployable until alignment is restored.

\subsection*{Stage 1 Contents Map}
\phantomsection
\addcontentsline{toc}{subsection}{Stage 1 Contents Map}

\begin{itemize}
\item 1.A Project Charter — Locks mission, scope boundaries, non-negotiables, definition-of-success posture aligned to Stage 0, decision hierarchy, roles and decision rights, and charter-level change-control posture.
\item 1.B Key Risks and Mitigations — Owned risk register with triggers or signals, impacts, likelihood, mitigations, and Owners; constrains gates and prevents unsafe escalation.
\item 1.C Assumptions and Dependencies — Converts prerequisites into validated dependencies with rationale, validation methods, due-by timing, Owners, and explicit fail actions enforcing substitution, delay, Hold, Regress, or Medical Hold.
\item 1.D Milestone Plan and Output Order — Locks stage build sequence and prerequisite gates; prevents out-of-order artifacts and contradiction-driven rework.
\item 1.E Standardized Templates — Mandatory template infrastructure, Overview Template and Phase Row Template, enforcing completeness, comparability, and auditability.
\item 1.F Governance Rules and Validity Doctrine — Formalizes validity tagging, gate outcomes, stop-rule integration, freeze windows, patch protocol, and admissibility rules restricting advancement decisions to Valid evidence only.
\end{itemize}

\end{document}